\SourceLeaf{198}
\Needspace{8\baselineskip}
\pdfbookmark[1]{Seeber, positive ternary quadratic forms. Review of July 9, 1831}{notice-1831-07-09}
\begin{center}
\rule{.86\linewidth}{.3pt}\par
Göttingische gelehrte Anzeigen.\quad July 9, 1831.\par
\rule{.86\linewidth}{.3pt}
\end{center}
\emph{Investigations into the Properties of Positive Ternary Quadratic Forms, by} \textsc{Ludwig August Seeber}, \emph{Doctor of Philosophy, Full Professor of Physics at the University in Freiburg. Freiburg im Breisgau} 1831. (248 pages in 4.)

Functions of two indeterminate quantities $x$ and $y$ of the form $axx+2bxy+cyy$, where $a,b,c$ denote specified integers, are known as \emph{quadratic forms}, or, when a further distinction is needed, \emph{binary quadratic forms}. They form one of the most interesting and richest subjects of higher arithmetic. The problems that arise first are these: to decide whether a given form of this kind contains another, $a'x'x'+2b'x'y'+c'y'y'$, that is, whether it can be transformed into that form by a substitution $x=\alpha x'+\beta y'$, $y=\gamma x'+\delta y'$, in which $\alpha,\beta,\gamma,\delta$ are integers; whether such a relation between two forms is mutual, in which case the forms are called equivalent; further, in both cases, to give all possible transformations of one into the other; and finally, to find all possible representations of a given integer by a given form with integer values of the indeterminate quantities. These problems are completely solved in the \emph{Disquisitiones Arithmeticae}, but constitute by far the smaller part of that work's section on quadratic forms. The finer investigations that follow require\textsuperscript{.} in part a preliminary treatment of a field one level higher and presenting much greater difficulties: the theory of similar functions of three indeterminate quantities $x,y,z$, which therefore have the form $axx+byy+czz+2a'yz+2b'xz+2c'xy$ and are called ternary quadratic forms. The solution of the principal problems concerning these ternary forms is developed in the work mentioned, though only as far as was necessary for the stated purpose. After an interval of thirty years, the author of the present work has now been the first to take up these investigations again and, for one principal kind of ternary form, namely the positive forms, has brought to completion what was left unfinished in the \emph{Disquisitiones }
\SourceLeaf{199}
\emph{Arithmeticae}. For those who have studied higher arithmetic more deeply, we could describe in a few words what is newly achieved in the present work. But to be understood by others too, we must allow ourselves somewhat more detail, and we do this all the more willingly because these investigations have an interest of their own outside higher arithmetic as well.

The properties of a binary form $axx+2bxy+cyy$ depend chiefly on the number $bb-ac$, which is therefore called the determinant of that form. Two equivalent forms always have equal determinants. But not all forms having a given determinant are therefore equivalent. Rather, such forms fall into a smaller or larger, but always finite, number of classes, with the forms belonging to one class equivalent to one another, and those belonging to different classes not equivalent. Forms whose determinant is positive can represent positive and negative numbers alike; forms with a negative determinant, however, can represent only numbers having the same sign as $a$ and $c$, so that positive and negative forms are distinguished here. The simplest forms in each class have definite criteria, are called reduced forms, and can be regarded as representatives of the whole class.

Similar relations concerning ternary forms are established in the \emph{Disquisitiones Arithmeticae}. The determinant of the ternary form
\[
axx+byy+czz+2a'yz+2b'xz+2c'xy
\]
is the number
\[
aa'a'+bb'b'+cc'c'-abc-2a'b'c'
\]
Here too, equality of determinants is necessary for two forms to be equivalent, but is not sufficient. Rather, all forms with a specified determinant fall into a finite number of classes, in each of which the simplest forms may be called reduced and may, as it were, represent all the others. The distinction between positive and negative forms, however, behaves differently here than for binary forms. For every given determinant, whether positive or negative, there are some forms that can represent
\SourceLeaf{200}
positive and negative numbers alike (indifferent forms), and others that can represent either only positive or only negative numbers (positive or negative forms). But positive forms exist only for negative determinants, and negative forms only for positive ones. It is evident, moreover, that a form's character as indifferent, positive or negative belongs at the same time to the whole class to which it belongs. The present work is confined to positive forms, whose determinants must therefore be negative. Everything that holds for these evidently transfers of itself to negative forms, whereas the indifferent forms, entirely excluded from the work, require a quite different treatment.

As already mentioned, the theory of ternary forms in the \emph{Disquisitiones Arithmeticae} was developed only as far as was needed for its purpose there. Consequently, the problem of deciding whether two given ternary forms are equivalent had not yet been solved in full generality. It was indeed shown there how to find, for any given form, an equivalent form of the simplest kind, and that there can be only finitely many such reduced forms for any given determinant. But since each class contains several such reduced forms, which do not in all cases turn out \emph{immediately} to be equivalent, a criterion was still lacking by which their equivalence or non-equivalence could be determined with certainty. The author of the present work has now met this need for positive forms completely and with exemplary thoroughness. His procedure is, however, presented somewhat differently from the account we have just given, and from how the matter would stand if the concept of reduced positive forms included only the essential conditions of greatest simplicity. For positive forms, these conditions are that the numbers $a,b,c$, which are positive by their nature, must be no smaller than, respectively, $b'$ or $c'$, $a'$ or $c'$, and $a'$ or $b'$, without regard to signs. Mr. \textsc{Seeber} has added further qualifications to the concept of reduced forms such that each class can contain only one form of this kind, and must contain one. In view of a beautiful theorem found by Mr. \textsc{Seeber} from observed cases, to be mentioned further below, we give here the principal conditions that Mr. S. has incorporated into the concept of reduced forms:
\SourceLeaf{201}
1) among the numbers $a',b',c'$, no two may have opposite signs; 2) without regard to sign, $2b'$ and $2c'$ must be no greater than $a$, further $a$ and $2a'$ no greater than $b$, and $b$ no greater than $c$; 3) when $a',b',c'$ are all negative, twice the sum of these numbers must be no greater than $a+b$. We may omit here the remaining qualifications that must be added for certain special cases.

The principal content of the work consists first in solving the problem of finding, for any given positive form, an equivalent form having the character of a reduced form according to the established definition; and then in the rigorous proof of the theorem that two non-identical reduced forms cannot be equivalent, or, what is the same thing, that each class contains only one reduced form. We must do full justice to the spirit of thoroughness with which these matters are carried through. And although we must regret the very great length of treatment involved, perhaps discouraging to some, since the solution of the problem takes 41 pages and the proof of the theorem 91 pages, we would by no means have this taken as a criticism. When a difficult problem or theorem is to be solved or proved, the first step, always to be acknowledged with due thanks, is that a solution or proof should be found at all. The question whether this could have been done more easily and simply remains idle so long as the possibility is not at the same time established by actually doing it. We therefore consider it premature to dwell on this question here. --- The remainder of the work chiefly contains equally thorough solutions of the following problems: to decide whether a given form contains another given form not equivalent to it; to find all possible transformations of a given form into a given equivalent form or one merely contained in it; and finally, for a given determinant, to give all possible classes of positive ternary forms.

We must also note that Mr. \textsc{Seeber} has expressed ternary forms somewhat differently from the \emph{Disquisitiones Arithmeticae}, where the coefficients of the products $yz$, $xz$, $xy$ were deliberately assumed to be even integers, whereas Mr. S. also allows odd ones, and therefore
\SourceLeaf{202}
denotes by $a',b',c'$ what was denoted above by $2a'$, $2b'$, $2c'$. The greater generality obtained in this way is evidently only apparent, or at least unnecessary, since everything that can be said about such forms with odd coefficients also follows of itself if their doubles are considered instead. We therefore cannot approve this change, which also entails some loss of simplicity. One consequence has been that what Mr. \textsc{Seeber} calls the determinant is always four times the number bearing that name in the \emph{Disquisitiones Arithmeticae}. In the present review we have retained the terminology of the \emph{Disquisitiones Arithmeticae}.

For the last-mentioned problem (to give all possible reduced forms for each given determinant), Mr. \textsc{Seeber} used a theorem to obtain bounds for the first three coefficients. According to this theorem, their product $abc$ cannot exceed three times the determinant. Mr. \textsc{Seeber} rigorously proves this theorem. In the preface, however, he remarks that among more than 600 cases he examined, he found not a single one in which that product exceeded twice the determinant. He therefore considers it highly probable that this tighter bound holds universally, but says that he did not succeed in finding a rigorous proof. Since this theorem, found by Mr. \textsc{Seeber} from observed cases, is both remarkable in itself and important for shortening the solution of the stated problem, we wish to add a very simple proof here, so that we too may contribute in this review to the further development of the theory. Two cases must be distinguished.

I. If none of the numbers $a',b',c'$ is negative, put
\[
\begin{aligned}
b-2a'&=d,&c-2b'&=e,&a-2c'&=f\\
c-2a'&=g,&a-2b'&=h,&b-2c'&=i
\end{aligned}
\]
The definition of reduced positive forms immediately implies that if
\[
axx+byy+czz+2a'yz+2b'xz+2c'xy
\]
is such a form, none of those six numbers is negative; it is also evident that $a,b,c$ are positive. If we now denote the form's (negative)
\SourceLeaf{203}
determinant by $-D$, we have the identity, readily checked by expansion,
\[
2D-abc=aa'd+bb'e+cc'f+a'hi+b'gi+c'gh+ghi
\]
in which none of the seven terms on the right can be negative, and consequently $abc$ is no greater than $2D$. The same conclusion follows in the same way from the identity
\[
2D-abc=aa'g+bb'h+cc'i+a'ef+b'df+c'de+def
\]

\Needspace{23\baselineskip}II. If none of the numbers $a',b',c'$ is positive, put
\[
\begin{aligned}
b+2a'&=d,&c+2b'&=e,&a+2c'&=f\\
c+2a'&=g,&a+2b'&=h,&b+2c'&=i
\end{aligned}
\]
\[
\begin{aligned}
b+c+2a'+2b'+2c'&=k\\
a+c+2a'+2b'+2c'&=l\\
a+b+2a'+2b'+2c'&=m
\end{aligned}
\]
and take the determinant of the form, as before, $=-D$. By the definition of reduced positive forms, none of the nine numbers $d,\allowbreak e,\allowbreak f,\allowbreak g,\allowbreak h,\allowbreak i,\allowbreak k,\allowbreak l,\allowbreak m$ can be negative. Thus the identity
\[
\begin{aligned}
6D-3abc={}&-aa'(d+2k)-bb'(e+2l)-cc'(f+2m)\\
&-a'hi-b'gi-c'gh+def+2ghi
\end{aligned}
\]
in which all terms on the right are positive or zero, since $a',b',c'$ are not positive but negative or zero, shows that $3abc$ cannot exceed $6D$, or $abc$ exceed $2D$. The same conclusion follows likewise from the identity
\[
\begin{aligned}
6D-3abc={}&-aa'(g+2k)-bb'(h+2l)-cc'(i+2m)\\
&-a'ef-b'df-c'de+2def+ghi
\end{aligned}
\]
Both equations are symmetric. If full symmetry is given up, the proof can be carried out with an even smaller number of terms, for example by the identity
\[
\begin{aligned}
8D-4abc={}&-2aa'(g+k)-2bb'(e+l)-4cc'm\\
&+(c+e)df+(c+g)hi
\end{aligned}
\]

\SourceLeaf{204}
We now wish to add a few remarks on the meaning of positive binary and ternary quadratic forms outside higher arithmetic. It is unnecessary to treat negative forms separately, while indifferent forms lie wholly outside this treatment.

The positive binary form $axx+2bxy+cyy$ represents, in general, the square of the distance between two unspecified points in a plane whose coordinates, relative to two axes inclined to one another at an angle with cosine $=\frac{b}{\surd ac}$, differ by $x\surd a$, $y\surd c$. If $x$ and $y$ are to mean only \emph{integers}, the form therefore refers to a system of points arranged in parallelograms, lying at the intersections of two systems of parallel lines. The lines in each system are equally spaced: the spacings in one, measured parallel to the lines in the second, are $=\surd a$, and the spacings in the other, measured parallel to the lines in the first, $=\surd c$; the inclination of the two systems is the one given above. The plane thus appears divided entirely into equal parallelograms whose vertices make up the system of points, with none of these points lying inside a parallelogram. The determinant taken with positive sign, thus $ac-bb$, is the square of the area of an elementary parallelogram. One and the same system of such points can be divided into parallelograms in infinitely many different ways and thus referred to equally many different forms. All these different forms are, however, what is technically called equivalent, and the area of an elementary parallelogram always remains the same. Two forms that are not equivalent, but of which one contains the other, refer to the same system of points: the first form to the whole system, the second to a part. Two forms that are technically called improperly equivalent (improprie aequivalentes) refer to two equal but oppositely placed systems of points, the plane being thought of as turned over, and so on.

In the same way, the positive ternary form
\[
axx+byy+czz+2a'yz+2b'xz+2c'xy
\]
represents, in general, the square of the distance between two unspecified points in space whose coordinates relative to three axes (1), (2), (3) give the differences $x\surd a$, $y\surd b$, $z\surd c$:
\SourceLeaf{205}
the cosines of the angles between axes (2) and (3), (1) and (3), and (1) and (2) are respectively $\frac{a'}{\surd bc}$, $\frac{b'}{\surd ac}$, $\frac{c'}{\surd ab}$. If $x,y,z$ are to mean only integers, the form refers to a system of points arranged in parallelepipeds, that is, points arising as the intersections of three systems of parallel, equally spaced planes. The whole space thus appears divided into equal parallelepipeds whose vertices make up that system of points, and the square of the volume of an elementary parallelepiped equals the determinant of the ternary form taken with positive sign\textsuperscript{.}. Equivalent forms represent one and the same system of points, referred only to other axes or fundamental planes. In the same way, all the other principal aspects of the theory of ternary forms find their geometric meaning here: one form being contained in another; the representation of a specified number or an indeterminate binary form by a ternary form; the theory of the adjoint ternary forms (formae adiunctae); the disappearance of the distinction between proper and improper equivalence; the nature of reduced forms; and so on. We must, however, confine ourselves to the indications above, especially since the present work, which considers ternary forms solely from a purely arithmetic standpoint, has given only indirect occasion for them. One will at least see from this what a rich field is opened here to investigations that have great theoretical interest in themselves and can also be used for a treatment of all relations among crystal forms that is both convenient and general. This is not the place to enter into the details of this application. We must nevertheless observe that, although $a,b,c,a',b',c'$ were originally assumed to denote integers, the greater part of the theory of ternary forms, and in particular what is required for that application, remains valid independently of this assumption. Indeed, \textsc{Haüy}'s specifications lead, for most kinds of crystal, to very simple integer values of the coefficients in the ternary forms referring to the corresponding arrangement of the system of points. But the more accurate later measurements of \textsc{Wollaston}, \textsc{Malus}, \textsc{Biot}, \textsc{Kupffer} and others contradict this and make it doubtful whether rational ratios of those coefficients agree with nature in every case. In any case, however, if one does not wish to drop from the theory the
\SourceLeaf{206}
restriction to integer values of the coefficients, integers can always be found that come as close as one wishes to the results of measurement, since what matters here is their ratios to one another, rather than their absolute values.

Finally, we wish to give the geometric meaning of \textsc{Seeber}'s theorem mentioned above. If a parallelepiped is such that none of its twelve edges (of which four at a time are equal) is greater than any of the twelve face diagonals (which are equal in pairs) or any of the four diagonals of the parallelepiped, then its volume multiplied by $\surd2$ is no smaller than the volume of a rectangular parallelepiped formed from the same edges.
\BookRule{.28\linewidth}
